\scanpage{7}
On en conclut la suite exacte
\[
0\longrightarrow E_0\otimes_{A_0}I\longrightarrow
\underset{\substack{\Vert\\E}}{H_D^0(A'/A,E')}\longrightarrow
\underset{\substack{\Vert\\E_0}}{H_D^0(A'_0/A_0,E'_0)}\longrightarrow0
\]
compte tenu que
\[
H^{\deleted{\unsure{0}}i}(C_{E_0\otimes_{A_0}I})=
\left\{\begin{aligned}
&E_0\otimes_{A_0}I &&\text{pour }i=0\\
&0 &&\text{pour }i=1\\
&&&(\text{et }=\ i\ne0)
\end{aligned}\right.
\]
Cela prouve d'abord que $E\longrightarrow E_0$ est \uuline{surjectif}, \deleted{donc $\operatorname{Im}(E\longrightarrow E')$ engendre le $A'$-module $E'$, puis} et de façon précise que

a) $\operatorname{Ker}(E\mathrel{\unsure{-|\longrightarrow}} E_0)=IE$, donc $E\longrightarrow E_0$ induit un \uline{isomorphisme} $E/IE=E\otimes_A A_0\xrightarrow{\sim}E_0$

b) $\underline{E/IE}\otimes_{A_0}I\longrightarrow IE$ un \uline{isomorphisme}, donc grâce à a) et $E_0$ plat sur $A_0$, $E$ est \uuline{plat} sur $A$

Considérons maintenant \deleted{le diagramme \unsure{comm.}} \inserted{l'hom}
\diagram{\begin{tikzcd}[ampersand replacement=\&]
E\otimes_A A'\arrow[r]\arrow[d,"\deleted{\mathrm{surj}}","\diagdown"{description,fill=none,inner sep=0pt}] \& E'\arrow[d,"\diagdown"{description,fill=none,inner sep=0pt}]\\
\deleted{E_0\otimes_{A_0}A'_0}\arrow[r,"\deleted{\sim}","\diagdown"{description,fill=none,inner sep=0pt}] \& \deleted{E'_0}
\end{tikzcd}}
c'est un hom de $A'$-modules \uline{plats} qui \deleted{on lit sur \unsure{le diagramme|les diagrammes} que le composé $E\otimes_A A'\longrightarrow E'\longrightarrow E'_0$ est surjectif, donc par Nakayama sur $A'$ (muni de $IA'$ nilpotent) $E\otimes_A A'$} \unsure{devient|donne} un isom. par réduction mod $IA'$, donc est un isomorphisme.
% archive sheet number: 6
